\scanpage{28}

On voit de $=$

\uline{Corollaire} 10\quad Conditions équivalentes

(i) $\mathcal F$ satisfait à $P_1$ ($P_2,P_3$)

(ii) On a $\underline H^i_Y(F)=0$ pour $i<1$ ($i<2$, $i<3$).

Reste le cas $P_3$, où on a posé

\[
\begin{aligned}
\underline H^2_Y(F)&=\text{faisceau associé du préfaisceau}\\
&\qquad V\rightsquigarrow \unsure{H^1|H1}(V\cap U,F).
\end{aligned}
\]

Or, $(\mathrm{i})\Rightarrow(\mathrm{ii})$ puisque $H^1(V\cap U,F)\simeq H^1(V,F)$ (\unsure{fonctoriellet|fonctoriellemnt} en $V$) et que le faisceau associé à $V\rightsquigarrow H^1(V,F)$ est trivial.

On a $(\mathrm{ii})\Rightarrow(\mathrm{i})$. \deleted{On \unsure{est|sera}} ramené en effet à prouver que tt $\unsure{\xi|\zeta|f}'\in\mathrm{Ob}\,\mathcal F(U)$ provient d'un $\xi\in\mathrm{Ob}\,\mathcal F(X)$. Or, \unsure{à} cause des propriétés d'unicité, c'est \uline{local} sur $Y$, et on procède alors comme dans la dém. de prop. 7, en \deleted{\unsure{tout}} raisonnant en un pt quelconque $x\in Y$.

% archive sheet number: 27
